\documentclass[aspectratio=169]{beamer}
%--- Configuración del Tema y Colores ---
\usetheme{Madrid}
\usecolortheme{whale}
\setbeamertemplate{navigation symbols}{}
\setbeamertemplate{caption}[numbered]

%--- Paquetes ---
\usepackage[utf8]{inputenc}
\usepackage[spanish]{babel}
\usepackage{graphicx}
\usepackage{booktabs}
\usepackage{tikz}
\usepackage{listings}
\usepackage{xcolor}
\usepackage{fontawesome5}
\usetikzlibrary{shapes.geometric, arrows, positioning}

%--- Configuración de Listings para Python ---
\definecolor{codegreen}{rgb}{0,0.6,0}
\definecolor{codegray}{rgb}{0.5,0.5,0.5}
\definecolor{codepurple}{rgb}{0.58,0,0.82}
\definecolor{backcolour}{rgb}{0.95,0.95,0.92}
\definecolor{promptcolor}{rgb}{0.0,0.4,0.6}

\lstdefinestyle{pythonstyle}{
    backgroundcolor=\color{backcolour},
    commentstyle=\color{codegreen},
    keywordstyle=\color{blue}\bfseries,
    numberstyle=\tiny\color{codegray},
    stringstyle=\color{codepurple},
    basicstyle=\ttfamily\scriptsize,
    breakatwhitespace=false,
    breaklines=true,
    captionpos=b,
    keepspaces=true,
    numbers=left,
    numbersep=5pt,
    showspaces=false,
    showstringspaces=false,
    showtabs=false,
    tabsize=2,
    language=Python,
    morekeywords={self, True, False, None, class, def, return, if, else, elif, for, while, in, and, or, not, lambda, async, await}
}

\lstdefinestyle{promptstyle}{
    backgroundcolor=\color{blue!5},
    basicstyle=\ttfamily\scriptsize\color{promptcolor},
    breaklines=true,
    frame=single,
    rulecolor=\color{blue!30},
}

\lstset{style=pythonstyle}

%--- Metadatos del Documento ---
\title[M503]{Programación Matemática 2}
\subtitle{Introducción al Vibe Coding\\Entornos Asistidos por IA}
\author[Luis Alfredo Alvarado]{Mgtr. Luis Alfredo Alvarado Rodríguez}
\institute[ECFM - USAC]{
    Escuela de Ciencias Físicas y Matemáticas \\
    Universidad de San Carlos de Guatemala
}
\date{Primer Semestre, 2026}

\begin{document}

%--- Diapositiva de Título ---
\begin{frame}
    \titlepage
\end{frame}

%--- Diapositiva: Tabla de Contenidos ---
\begin{frame}{Agenda de Hoy}
    \tableofcontents
\end{frame}

%===========================================================
% SECCIÓN 1: INTRODUCCIÓN AL VIBE CODING
%===========================================================
\section{¿Qué es el Vibe Coding?}

\begin{frame}{El Origen del Término}
    \begin{block}{Definición (Andrej Karpathy, 2025)}
        \textbf{``Vibe Coding''} es un estilo de programación donde el desarrollador se entrega a las vibes del código, acepta exponenciales, y olvida que el código siquiera existe.
    \end{block}
    \vspace{0.3cm}
    
    \begin{alertblock}{En la Práctica}
        Se describe completamente lo que se desea construir en lenguaje natural, y la IA genera el código. El rol del programador pasa de \textit{escribir} código a \textit{orquestar} y \textit{validar} código.
    \end{alertblock}
    \vspace{0.3cm}

\end{frame}

\begin{frame}{Evolución de la Programación}
    \centering
    \begin{tikzpicture}[node distance=1.8cm, auto]
        \tikzstyle{era} = [rectangle, rounded corners, minimum width=2.5cm, minimum height=1cm, text centered, draw=black, fill=blue!20]
        \tikzstyle{arrow} = [thick,->,>=stealth]
        
        \node (era1) [era] {1950s-70s\\Código Máquina};
        \node (era2) [era, right=0.8cm of era1] {1980s-90s\\Alto Nivel};
        \node (era3) [era, right=0.8cm of era2] {2000s-10s\\Frameworks};
        \node (era4) [era, right=0.8cm of era3, fill=green!30] {2020s+\\Vibe Coding};
        
        \draw [arrow] (era1) -- (era2);
        \draw [arrow] (era2) -- (era3);
        \draw [arrow] (era3) -- (era4);
        
        \node [below=0.3cm of era1, font=\tiny] {Assembly};
        \node [below=0.3cm of era2, font=\tiny] {C, Java, Python};
        \node [below=0.3cm of era3, font=\tiny] {React, Django};
        \node [below=0.3cm of era4, font=\tiny] {Copilot, Cursor};
    \end{tikzpicture}
    
    \vspace{0.5cm}
    \textbf{Tendencia:} Cada era \textbf{aumenta el nivel de abstracción} entre el programador y la máquina.
\end{frame}

\begin{frame}{El Nuevo Flujo de Trabajo}
    \begin{columns}[t]
        \column{0.48\textwidth}
        \begin{alertblock}{Programación Tradicional}
            \begin{enumerate}
                \item Pensar el algoritmo
                \item Escribir código línea por línea
                \item Debuggear errores de sintaxis
                \item Buscar en Stack Overflow
                \item Copiar y adaptar soluciones
            \end{enumerate}
        \end{alertblock}
        
        \column{0.48\textwidth}
        \begin{exampleblock}{Vibe Coding}
            \begin{enumerate}
                \item Describir la intención en lenguaje natural
                \item La IA genera el código
                \item Revisar y validar la solución
                \item Iterar con nuevos prompts
                \item Refinar hasta obtener el resultado
            \end{enumerate}
        \end{exampleblock}
    \end{columns}
    
    \vspace{0.5cm}
    \centering
    \textbf{Cambio de rol:} De \textit{escritor} de código a \textit{director} de código.
\end{frame}

%===========================================================
% SECCIÓN 2: HERRAMIENTAS PRINCIPALES
%===========================================================
\section{Herramientas de Programación Asistida}

\begin{frame}{Ecosistema de Herramientas (2026)}
    \begin{table}[]
        \centering
        \renewcommand{\arraystretch}{1.3}
        \scriptsize
        \begin{tabular}{p{2.5cm} p{3cm} p{4cm} p{3cm}}
            \toprule
            \textbf{Herramienta} & \textbf{Empresa} & \textbf{Características} & \textbf{Modelo Base} \\
            \midrule
            GitHub Copilot & Microsoft/GitHub & Integración VS Code, autocompletado & GPT-4 / Codex \\
            Cursor & Cursor Inc. & IDE completo, chat integrado & Claude / GPT-4 \\
            Claude Code & Anthropic & Terminal, agéntico & Claude 3.5/4 \\
            Windsurf & Codeium & Flujo agéntico, gratuito & Modelos propios \\
            Amazon Q & AWS & Integración AWS & Titan / Claude \\
            Gemini Code & Google & Integración Google Cloud & Gemini \\
            \bottomrule
        \end{tabular}
    \end{table}
    
    \vspace{0.3cm}
    \centering
    \textbf{En este curso:} Se utilizará principalmente \textbf{Cursor} y \textbf{GitHub Copilot}.
\end{frame}

\begin{frame}{GitHub Copilot}
    \begin{columns}
        \column{0.6\textwidth}
        \textbf{Características principales:}
        \begin{itemize}
            \item Autocompletado inteligente en tiempo real
            \item Sugerencias basadas en contexto del archivo
            \item Generación de funciones completas desde comentarios
            \item Soporte para múltiples lenguajes
            \item Chat integrado (Copilot Chat)
        \end{itemize}
        
        \vspace{0.3cm}
        \textbf{Integración:}
        \begin{itemize}
            \item VS Code, Visual Studio, JetBrains
            \item Neovim, Vim
        \end{itemize}
        
        \column{0.4\textwidth}
        \begin{block}{Precio (2026)}
            \begin{itemize}
                \item Individual: \$10/mes
                \item Business: \$19/mes
                \item \textbf{Estudiantes: Gratis}
            \end{itemize}
        \end{block}
        
        \vspace{0.3cm}
        \centering
        \faGithub \ \texttt{github.com/features/copilot}
    \end{columns}
\end{frame}

\begin{frame}{Cursor: El IDE del Futuro}
    \begin{columns}
        \column{0.6\textwidth}
        \textbf{Características principales:}
        \begin{itemize}
            \item IDE completo basado en VS Code
            \item Chat con contexto del proyecto completo
            \item Modo ``Composer'' para cambios multi-archivo
            \item Comando Cmd+K para edición inline
            \item Integración con múltiples LLMs
            \item Indexación semántica del codebase
        \end{itemize}
        
        \vspace{0.3cm}
        \textbf{Modos de interacción:}
        \begin{itemize}
            \item \textbf{Tab:} Autocompletado
            \item \textbf{Chat:} Conversación contextual
            \item \textbf{Composer:} Edición multi-archivo
        \end{itemize}
        
        \column{0.4\textwidth}
        \begin{block}{Precio (2026)}
            \begin{itemize}
                \item Hobby: Gratis (limitado)
                \item Pro: \$20/mes
                \item Business: \$40/mes
            \end{itemize}
        \end{block}
        
        \vspace{0.3cm}
        \centering
        \faDesktop \ \texttt{cursor.com}
    \end{columns}
\end{frame}

\begin{frame}{Comparación: Copilot vs Cursor}
    \begin{table}[]
        \centering
        \renewcommand{\arraystretch}{1.4}
        \begin{tabular}{p{4cm} p{4.5cm} p{4.5cm}}
            \toprule
            \textbf{Aspecto} & \textbf{GitHub Copilot} & \textbf{Cursor} \\
            \midrule
            Tipo & Plugin/Extensión & IDE Completo \\
            Contexto & Archivo actual & Proyecto completo \\
            Edición multi-archivo & Limitada & Nativa (Composer) \\
            Modelos & GPT-4 (fijo) & Claude, GPT-4, etc. \\
            Curva de aprendizaje & Baja & Media \\
            Ideal para & Autocompletado rápido & Proyectos complejos \\
            \bottomrule
        \end{tabular}
    \end{table}
    
    \vspace{0.3cm}
    \centering
    \textbf{Recomendación:} Comenzar con Copilot, migrar a Cursor para proyectos mayores.
\end{frame}

%===========================================================
% SECCIÓN 3: CÓMO FUNCIONAN
%===========================================================
\section{¿Cómo Funcionan los Asistentes de Código?}

\begin{frame}{Arquitectura General}
    \centering
    \begin{tikzpicture}[node distance=1.5cm, auto]
        \tikzstyle{block} = [rectangle, rounded corners, minimum width=2.5cm, minimum height=1cm, text centered, draw=black, fill=blue!20]
        \tikzstyle{arrow} = [thick,->,>=stealth]
        
        \node (input) [block] {Código + Prompt};
        \node (tokenizer) [block, right=1cm of input] {Tokenizador};
        \node (llm) [block, right=1cm of tokenizer, fill=green!30] {LLM};
        \node (output) [block, right=1cm of llm] {Código Generado};
        
        \draw [arrow] (input) -- (tokenizer);
        \draw [arrow] (tokenizer) -- (llm);
        \draw [arrow] (llm) -- (output);
        
        \node [below=0.5cm of input, font=\scriptsize, text width=2.5cm, align=center] {Contexto del archivo, comentarios, instrucciones};
        \node [below=0.5cm of tokenizer, font=\scriptsize, text width=2.5cm, align=center] {Convierte texto a tokens numéricos};
        \node [below=0.5cm of llm, font=\scriptsize, text width=2.5cm, align=center] {GPT-4, Claude, etc.};
        \node [below=0.5cm of output, font=\scriptsize, text width=2.5cm, align=center] {Sugerencias, funciones, explicaciones};
    \end{tikzpicture}
    
    \vspace{0.5cm}
    El asistente \textbf{predice} el código más probable dado el contexto proporcionado.
\end{frame}

\begin{frame}{El Contexto lo es Todo}
    \begin{block}{¿Qué información utiliza el asistente?}
        \begin{itemize}
            \item Código anterior y posterior al cursor
            \item Nombres de archivos y estructura del proyecto
            \item Comentarios y docstrings
            \item Imports y dependencias
            \item Historial de la conversación (en modo chat)
        \end{itemize}
    \end{block}
    
    \vspace{0.3cm}
    \begin{alertblock}{Implicación Importante}
        \textbf{Código bien documentado = Mejores sugerencias de IA}
        
        Los nombres descriptivos, comentarios claros y estructura organizada mejoran dramáticamente la calidad del código generado.
    \end{alertblock}
\end{frame}

\begin{frame}{Limitaciones Técnicas}
    \begin{columns}[t]
        \column{0.48\textwidth}
        \begin{alertblock}{Ventana de Contexto}
            \begin{itemize}
                \item Los LLMs tienen límite de tokens
                \item GPT-4: ~128K tokens
                \item Claude 3: ~200K tokens
                \item Proyectos grandes exceden el límite
            \end{itemize}
        \end{alertblock}
        
        \column{0.48\textwidth}
        \begin{alertblock}{Alucinaciones}
            \begin{itemize}
                \item La IA puede inventar APIs inexistentes
                \item Puede generar código que ``parece'' correcto
                \item No tiene acceso a documentación actualizada
                \item \textbf{Siempre verificar el código generado}
            \end{itemize}
        \end{alertblock}
    \end{columns}
    
    \vspace{0.5cm}
    \centering
    \textit{``La IA es un copiloto, no un piloto automático.''}
\end{frame}

%===========================================================
% SECCIÓN 4: FLUJO DE TRABAJO PRÁCTICO
%===========================================================
\section{Flujo de Trabajo con Vibe Coding}

\begin{frame}{El Ciclo del Vibe Coding}
    \centering
    \begin{tikzpicture}[node distance=2cm, auto]
        \tikzstyle{step} = [circle, minimum size=1.5cm, text centered, draw=black, fill=blue!20, font=\scriptsize]
        \tikzstyle{arrow} = [thick,->,>=stealth]
        
        \node (describe) [step, fill=green!30] {1. Describir};
        \node (generate) [step, right=1.5cm of describe] {2. Generar};
        \node (review) [step, right=1.5cm of generate] {3. Revisar};
        \node (test) [step, below=1.5cm of review] {4. Probar};
        \node (iterate) [step, left=1.5cm of test] {5. Iterar};
        
        \draw [arrow] (describe) -- (generate);
        \draw [arrow] (generate) -- (review);
        \draw [arrow] (review) -- (test);
        \draw [arrow] (test) -- (iterate);
        \draw [arrow] (iterate) -- (describe);
        
        \node [right=0.3cm of describe, font=\tiny, text width=2cm] {Escribir prompt claro};
        \node [right=0.3cm of review, font=\tiny, text width=2cm] {Verificar lógica};
        \node [below=0.3cm of test, font=\tiny, text width=2cm] {Ejecutar y validar};
    \end{tikzpicture}
    
    \vspace{0.5cm}
    \textbf{Clave:} La iteración rápida. No se espera perfección en el primer intento.
\end{frame}

\begin{frame}{Anatomía de un Buen Prompt}
    \begin{exampleblock}{Estructura Recomendada}
        \begin{enumerate}
            \item \textbf{Contexto:} ¿Qué se está construyendo?
            \item \textbf{Tarea específica:} ¿Qué debe hacer esta función/clase?
            \item \textbf{Restricciones:} ¿Hay requisitos especiales?
            \item \textbf{Formato:} ¿Cómo debe verse el resultado?
        \end{enumerate}
    \end{exampleblock}
    
    \vspace{0.3cm}
    \begin{columns}
        \column{0.48\textwidth}
        \begin{alertblock}{Prompt Malo \faTimesCircle}
            ``Haz una función para ordenar''
        \end{alertblock}
        
        \column{0.48\textwidth}
        \begin{exampleblock}{Prompt Bueno \faCheckCircle}
            ``Crear una función que ordene una lista de diccionarios por el campo 'fecha' en orden descendente. Usar el módulo datetime.''
        \end{exampleblock}
    \end{columns}
\end{frame}

\begin{frame}[fragile]{Ejemplo: Prompt $\rightarrow$ Código}
    \begin{lstlisting}[style=promptstyle, title=Prompt en Cursor/Copilot]
Crear una clase CuentaBancaria con:
- Atributos: titular (str), saldo (float, privado)
- Metodos: depositar, retirar (con validacion), obtener_saldo
- El saldo no puede ser negativo
- Incluir docstrings en espanol
    \end{lstlisting}
    
    \vspace{0.2cm}
    \textbf{La IA genera:}
    \begin{lstlisting}[basicstyle=\ttfamily\tiny]
class CuentaBancaria:
    """Representa una cuenta bancaria con operaciones basicas."""
    
    def __init__(self, titular: str, saldo_inicial: float = 0.0):
        """Inicializa la cuenta con un titular y saldo inicial."""
        self.titular = titular
        self._saldo = max(0.0, saldo_inicial)
    
    def depositar(self, monto: float) -> bool:
        """Deposita un monto positivo en la cuenta."""
        if monto > 0:
            self._saldo += monto
            return True
        return False
    \end{lstlisting}
\end{frame}

\begin{frame}{Técnicas Avanzadas de Prompting}
    \begin{enumerate}
        \setlength\itemsep{0.8em}
        \item \textbf{Few-shot prompting:} Proporcionar ejemplos del formato deseado antes de la solicitud.
        
        \item \textbf{Chain of thought:} Pedir que explique el razonamiento paso a paso.
        
        \item \textbf{Iteración incremental:} Construir el código en partes pequeñas, no todo de una vez.
        
        \item \textbf{Contexto explícito:} Incluir información relevante del proyecto en el prompt.
        
        \item \textbf{Restricciones negativas:} Especificar lo que \textbf{no} se desea (``sin usar librerías externas'').
    \end{enumerate}
\end{frame}

%===========================================================
% SECCIÓN 5: DEMOSTRACIÓN PRÁCTICA
%===========================================================
\section{Demostración Práctica}

\begin{frame}[fragile]{Demo 1: Autocompletado Básico}
    \textbf{Se escribe un comentario descriptivo:}
    \begin{lstlisting}
# Funcion que calcula el factorial de un numero usando recursion
    \end{lstlisting}
    
    \vspace{0.3cm}
    \textbf{Copilot/Cursor sugiere automáticamente:}
    \begin{lstlisting}
# Funcion que calcula el factorial de un numero usando recursion
def factorial(n: int) -> int:
    """Calcula el factorial de n de forma recursiva."""
    if n <= 1:
        return 1
    return n * factorial(n - 1)
    \end{lstlisting}
    
    \vspace{0.3cm}
    \textbf{Acción:} Presionar \texttt{Tab} para aceptar la sugerencia.
\end{frame}

\begin{frame}[fragile]{Demo 2: Generación desde Docstring}
    \textbf{Se escribe solo la firma y documentación:}
    \begin{lstlisting}
def fibonacci(n: int) -> list[int]:
    """
    Genera los primeros n numeros de la secuencia Fibonacci.
    
    Args:
        n: Cantidad de numeros a generar
        
    Returns:
        Lista con los primeros n numeros Fibonacci
        
    Example:
        >>> fibonacci(7)
        [0, 1, 1, 2, 3, 5, 8]
    """
    \end{lstlisting}
    
    \vspace{0.2cm}
    \textbf{La IA completa la implementación basándose en la documentación.}
\end{frame}

\begin{frame}[fragile]{Demo 3: Chat para Refactorización}
    \begin{lstlisting}[style=promptstyle, title=Prompt en el Chat]
Refactorizar este codigo para que sea mas eficiente y use list comprehension:

def obtener_pares(numeros):
    resultado = []
    for n in numeros:
        if n % 2 == 0:
            resultado.append(n)
    return resultado
    \end{lstlisting}
    
    \vspace{0.2cm}
    \textbf{Respuesta de la IA:}
    \begin{lstlisting}
def obtener_pares(numeros: list[int]) -> list[int]:
    """Retorna los numeros pares de la lista usando list comprehension."""
    return [n for n in numeros if n % 2 == 0]
    \end{lstlisting}
\end{frame}

\begin{frame}[fragile]{Demo 4: Generación de Tests}
    \begin{lstlisting}[style=promptstyle, title=Prompt]
Generar tests unitarios con pytest para la clase CuentaBancaria.
Cubrir casos: deposito valido, deposito invalido, retiro valido, 
retiro que excede saldo, saldo inicial negativo.
    \end{lstlisting}
    
    \vspace{0.2cm}
    \textbf{La IA genera:}
    \begin{lstlisting}[basicstyle=\ttfamily\tiny]
import pytest
from cuenta_bancaria import CuentaBancaria

class TestCuentaBancaria:
    def test_deposito_valido(self):
        cuenta = CuentaBancaria("Juan", 100)
        assert cuenta.depositar(50) == True
        assert cuenta.obtener_saldo() == 150
    
    def test_deposito_invalido(self):
        cuenta = CuentaBancaria("Juan", 100)
        assert cuenta.depositar(-50) == False
        assert cuenta.obtener_saldo() == 100
    \end{lstlisting}
\end{frame}

\begin{frame}[fragile]{Demo 5: Composer (Multi-archivo)}
    \begin{lstlisting}[style=promptstyle, title=Prompt en Cursor Composer]
Crear una API REST basica con FastAPI para gestionar tareas (TODO):
- Modelo: Tarea con id, titulo, descripcion, completada
- Endpoints: GET /tareas, POST /tareas, PUT /tareas/{id}, DELETE /tareas/{id}
- Usar Pydantic para validacion
- Crear archivos separados: models.py, routes.py, main.py
    \end{lstlisting}
    
    \vspace{0.3cm}
    \textbf{Cursor genera automáticamente los tres archivos con la estructura completa.}
    
    \vspace{0.3cm}
    \begin{exampleblock}{Ventaja del Composer}
        Se mantiene la coherencia entre archivos: imports correctos, referencias cruzadas, y convenciones consistentes.
    \end{exampleblock}
\end{frame}

%===========================================================
% SECCIÓN 6: BUENAS PRÁCTICAS
%===========================================================
\section{Buenas Prácticas y Consideraciones}

\begin{frame}{Las 7 Reglas del Vibe Coding}
    \begin{enumerate}
        \setlength\itemsep{0.6em}
        \item \textbf{Siempre revisar:} Nunca aceptar código sin entenderlo.
        \item \textbf{Probar todo:} Ejecutar y validar antes de integrar.
        \item \textbf{Iterar rápido:} Pequeños cambios, feedback inmediato.
        \item \textbf{Ser específico:} Prompts vagos = código vago.
        \item \textbf{Mantener contexto:} Nombres descriptivos mejoran las sugerencias.
        \item \textbf{Conocer los límites:} La IA no reemplaza el conocimiento fundamental.
        \item \textbf{Documentar decisiones:} Explicar por qué, no solo qué.
    \end{enumerate}
\end{frame}

\begin{frame}{Cuándo NO Usar Vibe Coding Ciegamente}
    \begin{alertblock}{Situaciones de Alto Riesgo}
        \begin{itemize}
            \item Código de seguridad (autenticación, criptografía)
            \item Sistemas críticos (médicos, financieros, aeronáuticos)
            \item Manejo de datos sensibles (PII, HIPAA, GDPR)
            \item Algoritmos con requisitos de rendimiento estrictos
        \end{itemize}
    \end{alertblock}
    
    \vspace{0.3cm}
    \begin{exampleblock}{En Estos Casos}
        \begin{itemize}
            \item Usar la IA como \textbf{punto de partida}, no como solución final
            \item Revisión exhaustiva por expertos
            \item Testing riguroso y auditorías de código
        \end{itemize}
    \end{exampleblock}
\end{frame}

\begin{frame}{Consideraciones Éticas y Legales}
    \begin{columns}[t]
        \column{0.48\textwidth}
        \begin{block}{Propiedad Intelectual}
            \begin{itemize}
                \item ¿De quién es el código generado?
                \item Los LLMs se entrenan con código público
                \item Posible reproducción de código con licencia
                \item Verificar licencias de dependencias
            \end{itemize}
        \end{block}
        
        \column{0.48\textwidth}
        \begin{block}{Privacidad}
            \begin{itemize}
                \item El código se envía a servidores externos
                \item No usar con código propietario sensible
                \item Revisar políticas de retención de datos
                \item Opciones on-premise para empresas
            \end{itemize}
        \end{block}
    \end{columns}
    
    \vspace{0.5cm}
    \centering
    \textbf{Recomendación:} Revisar las políticas de uso de cada herramienta.
\end{frame}

\begin{frame}{El Futuro del Desarrollador}
    \begin{exampleblock}{Habilidades que Aumentan de Valor}
        \begin{itemize}
            \item Pensamiento arquitectónico y diseño de sistemas
            \item Comunicación clara (para mejores prompts)
            \item Revisión crítica de código
            \item Comprensión profunda de fundamentos
            \item Creatividad y resolución de problemas complejos
        \end{itemize}
    \end{exampleblock}
    
    \vspace{0.3cm}
    \begin{alertblock}{Habilidades que Disminuyen de Valor}
        \begin{itemize}
            \item Memorización de sintaxis
            \item Tareas repetitivas de codificación
            \item Búsqueda manual de soluciones comunes
        \end{itemize}
    \end{alertblock}
\end{frame}

%===========================================================
% SECCIÓN 7: CONFIGURACIÓN DEL ENTORNO
%===========================================================
\section{Configuración del Entorno}

\begin{frame}{Instalación de Cursor}
    \textbf{Pasos para configurar Cursor:}
    \begin{enumerate}
        \setlength\itemsep{0.5em}
        \item Descargar desde \texttt{cursor.com}
        \item Instalar (disponible para Windows, Mac, Linux)
        \item Crear cuenta o iniciar sesión
        \item Importar configuración de VS Code (opcional)
        \item Seleccionar modelo preferido (Claude recomendado)
    \end{enumerate}
    
    \vspace{0.3cm}
    \begin{block}{Atajos Esenciales}
        \begin{itemize}
            \item \textbf{Cmd/Ctrl + K:} Edición inline con prompt
            \item \textbf{Cmd/Ctrl + L:} Abrir chat
            \item \textbf{Cmd/Ctrl + I:} Abrir Composer
            \item \textbf{Tab:} Aceptar sugerencia
            \item \textbf{Esc:} Rechazar sugerencia
        \end{itemize}
    \end{block}
\end{frame}

\begin{frame}{Instalación de GitHub Copilot}
    \textbf{Pasos para configurar Copilot en VS Code:}
    \begin{enumerate}
        \setlength\itemsep{0.5em}
        \item Abrir VS Code
        \item Ir a Extensions (Cmd/Ctrl + Shift + X)
        \item Buscar ``GitHub Copilot''
        \item Instalar la extensión oficial
        \item Autenticarse con cuenta de GitHub
        \item (Estudiantes) Verificar GitHub Education para acceso gratuito
    \end{enumerate}
    
    \vspace{0.3cm}
    \begin{exampleblock}{Verificar GitHub Education}
        \begin{enumerate}
            \item Ir a \texttt{education.github.com}
            \item Registrarse con correo institucional (@usac.edu.gt)
            \item Subir comprobante de estudiante
            \item Esperar aprobación (24-72 horas)
        \end{enumerate}
    \end{exampleblock}
\end{frame}

%===========================================================
% SECCIÓN 8: EJERCICIOS
%===========================================================
\section{Ejercicios Propuestos}

\begin{frame}{Ejercicios para Practicar}
    \textbf{Usar Cursor o Copilot para los siguientes ejercicios:}
    \vspace{0.3cm}
    
    \begin{enumerate}
        \setlength\itemsep{0.7em}
        \item Generar una función que valide si un string es un palíndromo (ignorando espacios y mayúsculas).
        
        \item Crear una clase \texttt{Calculadora} con operaciones básicas y un historial de operaciones.
        
        \item Generar una función que convierta números romanos a enteros y viceversa.
        
        \item Crear un script que lea un archivo CSV y genere estadísticas básicas (media, mediana, moda).
        
        \item \textbf{(Bonus)} Usar el Composer de Cursor para generar una aplicación CLI de lista de tareas con persistencia en JSON.
    \end{enumerate}
    
    \vspace{0.3cm}
    \textbf{Importante:} Documentar los prompts utilizados y las iteraciones necesarias.
\end{frame}

%--- Diapositiva Final ---
\begin{frame}
    \centering
    \Huge \textbf{¿Preguntas?}
    
    \vspace{0.8cm}
    \large
    \textbf{Próxima clase:} Iteración Rápida y Revisión de Código con IA
    
    \vspace{0.5cm}
    \normalsize
    Mgtr. Luis Alfredo Alvarado Rodríguez\\
    \small Escuela de Ciencias Físicas y Matemáticas
    
    \vspace{0.5cm}
    \footnotesize
    \textit{``La mejor manera de predecir el futuro es inventarlo.''} — Alan Kay
\end{frame}

\end{document}